\chapter{自动化}

% === 源文件: automation/auth-monitoring.md ===
\section{认证监控}


OpenClaw 通过 \texttt{openclaw models status} 提供 OAuth 过期健康状态。请使用该命令进行自动化和告警；脚本是为手机工作流程提供的可选附加功能。

\subsection{推荐方式：CLI 检查（可移植）}


\begin{lstlisting}[language=bash]
openclaw models status --check
\end{lstlisting}


退出码：

\begin{itemize}
  \item \texttt{0}：正常
  \item \texttt{1}：凭证过期或缺失
  \item \texttt{2}：即将过期（24 小时内）
\end{itemize}


此方式适用于 cron/systemd，无需额外脚本。

\subsection{可选脚本（运维 / 手机工作流程）}


这些脚本位于 \texttt{scripts/} 目录下，属于\textbf{可选}内容。它们假定你可以通过 SSH 访问 Gateway 网关主机，并针对 systemd + Termux 进行了调优。

\begin{itemize}
  \item \texttt{scripts/claude-auth-status.sh} 现在使用 \texttt{openclaw models status --json} 作为数据来源（如果 CLI 不可用则回退到直接读取文件），因此请确保 \texttt{openclaw} 在定时器的 \texttt{PATH} 中。
  \item \texttt{scripts/auth-monitor.sh}：cron/systemd 定时器目标；发送告警（ntfy 或手机）。
  \item \texttt{scripts/systemd/openclaw-auth-monitor.\{service,timer\}}：systemd 用户定时器。
  \item \texttt{scripts/claude-auth-status.sh}：Claude Code + OpenClaw 认证检查器（完整/json/简洁模式）。
  \item \texttt{scripts/mobile-reauth.sh}：通过 SSH 引导的重新认证流程。
  \item \texttt{scripts/termux-quick-auth.sh}：一键小部件状态查看 + 打开认证 URL。
  \item \texttt{scripts/termux-auth-widget.sh}：完整的引导式小部件流程。
  \item \texttt{scripts/termux-sync-widget.sh}：同步 Claude Code 凭证 → OpenClaw。
\end{itemize}


如果你不需要手机自动化或 systemd 定时器，可以跳过这些脚本。


\clearpage

% === 源文件: automation/cron-jobs.md ===
\section{定时任务（Gateway网关调度器）}


\begin{quote}
\textbf{定时任务还是心跳？} 请参阅\href{/automation/cron-vs-heartbeat}{定时任务与心跳对比}了解何时使用哪种方式。
\end{quote}


定时任务是 Gateway网关内置的调度器。它持久化任务、在合适的时间唤醒智能体，并可选择将输出发送回聊天。

如果你想要 \_"每天早上运行"\_ 或 \_"20 分钟后提醒智能体"\_，定时任务就是对应的机制。

\subsection{简要概述}


\begin{itemize}
  \item 定时任务运行在 \textbf{Gateway网关内部}（而非模型内部）。
  \item 任务持久化存储在 \texttt{\textasciitilde{}/.openclaw/cron/} 下，因此重启不会丢失计划。
  \item 两种执行方式：
  \item \textbf{主会话}：入队一个系统事件，然后在下一次心跳时运行。
  \item \textbf{隔离式}：在 \texttt{cron:} 中运行专用智能体轮次，可投递摘要（默认 announce）或不投递。
  \item 唤醒是一等功能：任务可以请求"立即唤醒"或"下次心跳时"。
\end{itemize}


\subsection{快速开始（可操作）}


创建一个一次性提醒，验证其存在，然后立即运行：

\begin{lstlisting}[language=bash]
openclaw cron add \
  --name "Reminder" \
  --at "2026-02-01T16:00:00Z" \
  --session main \
  --system-event "Reminder: check the cron docs draft" \
  --wake now \
  --delete-after-run

openclaw cron list
openclaw cron run <job-id> --force
openclaw cron runs --id <job-id>
\end{lstlisting}


调度一个带投递功能的周期性隔离任务：

\begin{lstlisting}[language=bash]
openclaw cron add \
  --name "Morning brief" \
  --cron "0 7 * * *" \
  --tz "America/Los_Angeles" \
  --session isolated \
  --message "Summarize overnight updates." \
  --announce \
  --channel slack \
  --to "channel:C1234567890"
\end{lstlisting}


\subsection{工具调用等价形式（Gateway网关定时任务工具）}


有关规范的 JSON 结构和示例，请参阅\href{/automation/cron-jobs\#json-schema-for-tool-calls}{工具调用的 JSON 模式}。

\subsection{定时任务的存储位置}


定时任务默认持久化存储在 Gateway网关主机的 \texttt{\textasciitilde{}/.openclaw/cron/jobs.json} 中。Gateway网关将文件加载到内存中，并在更改时写回，因此仅在 Gateway网关停止时手动编辑才是安全的。请优先使用 \texttt{openclaw cron add/edit} 或定时任务工具调用 API 进行更改。

\subsection{新手友好概述}


将定时任务理解为：\textbf{何时}运行 + \textbf{做什么}。

\begin{enumerate}
  \item \textbf{选择调度计划}
\end{enumerate}

\begin{itemize}
  \item 一次性提醒 → \texttt{schedule.kind = "at"}（CLI：\texttt{--at}）
  \item 重复任务 → \texttt{schedule.kind = "every"} 或 \texttt{schedule.kind = "cron"}
  \item 如果你的 ISO 时间戳省略了时区，将被视为 \textbf{UTC}。
\end{itemize}


\begin{enumerate}
  \item \textbf{选择运行位置}
\end{enumerate}

\begin{itemize}
  \item \texttt{sessionTarget: "main"} → 在下一次心跳时使用主会话上下文运行。
  \item \texttt{sessionTarget: "isolated"} → 在 \texttt{cron:} 中运行专用智能体轮次。
\end{itemize}


\begin{enumerate}
  \item \textbf{选择负载}
\end{enumerate}

\begin{itemize}
  \item 主会话 → \texttt{payload.kind = "systemEvent"}
  \item 隔离会话 → \texttt{payload.kind = "agentTurn"}
\end{itemize}


可选：一次性任务（\texttt{schedule.kind = "at"}）默认会在成功运行后删除。设置
\texttt{deleteAfterRun: false} 可保留它（成功后会禁用）。

\subsection{概念}


\subsubsection{任务}


定时任务是一条存储记录，包含：

\begin{itemize}
  \item 一个\textbf{调度计划}（何时运行），
  \item 一个\textbf{负载}（做什么），
  \item 可选的\textbf{投递}（输出发送到哪里）。
  \item 可选的\textbf{智能体绑定}（\texttt{agentId}）：在指定智能体下运行任务；如果缺失或未知，Gateway网关会回退到默认智能体。
\end{itemize}


任务通过稳定的 \texttt{jobId} 标识（用于 CLI/Gateway网关 API）。
在智能体工具调用中，\texttt{jobId} 是规范字段；旧版 \texttt{id} 仍可兼容使用。
一次性任务默认会在成功运行后自动删除；设置 \texttt{deleteAfterRun: false} 可保留它。

\subsubsection{调度计划}


定时任务支持三种调度类型：

\begin{itemize}
  \item \texttt{at}：一次性时间戳（ISO 8601 字符串）。
  \item \texttt{every}：固定间隔（毫秒）。
  \item \texttt{cron}：5 字段 cron 表达式，可选 IANA 时区。
\end{itemize}


Cron 表达式使用 \texttt{croner}。如果省略时区，将使用 Gateway网关主机的本地时区。

\subsubsection{主会话与隔离式执行}


\paragraph{主会话任务（系统事件）}


主会话任务入队一个系统事件，并可选择唤醒心跳运行器。它们必须使用 \texttt{payload.kind = "systemEvent"}。

\begin{itemize}
  \item \texttt{wakeMode: "next-heartbeat"}（默认）：事件等待下一次计划心跳。
  \item \texttt{wakeMode: "now"}：事件触发立即心跳运行。
\end{itemize}


当你需要正常的心跳提示 + 主会话上下文时，这是最佳选择。参见\href{/gateway/heartbeat}{心跳}。

\paragraph{隔离任务（专用定时会话）}


隔离任务在会话 \texttt{cron:} 中运行专用智能体轮次。

关键行为：

\begin{itemize}
  \item 提示以 \texttt{[cron: <任务名称>]} 为前缀，便于追踪。
  \item 每次运行都会启动一个\textbf{全新的会话 ID}（不继承之前的对话）。
  \item 如果未指定 \texttt{delivery}，隔离任务会默认以“announce”方式投递摘要。
  \item \texttt{delivery.mode} 可选 \texttt{announce}（投递摘要）或 \texttt{none}（内部运行）。
\end{itemize}


对于嘈杂、频繁或"后台杂务"类任务，使用隔离任务可以避免污染你的主聊天记录。

\subsubsection{负载结构（运行内容）}


支持两种负载类型：

\begin{itemize}
  \item \texttt{systemEvent}：仅限主会话，通过心跳提示路由。
  \item \texttt{agentTurn}：仅限隔离会话，运行专用智能体轮次。
\end{itemize}


常用 \texttt{agentTurn} 字段：

\begin{itemize}
  \item \texttt{message}：必填文本提示。
  \item \texttt{model} / \texttt{thinking}：可选覆盖（见下文）。
  \item \texttt{timeoutSeconds}：可选超时覆盖。
\end{itemize}


\subsubsection{模型和思维覆盖}


隔离任务（\texttt{agentTurn}）可以覆盖模型和思维级别：

\begin{itemize}
  \item \texttt{model}：提供商/模型字符串（例如 \texttt{anthropic/claude-sonnet-4-20250514}）或别名（例如 \texttt{opus}）
  \item \texttt{thinking}：思维级别（\texttt{off}、\texttt{minimal}、\texttt{low}、\texttt{medium}、\texttt{high}、\texttt{xhigh}；仅限 GPT-5.2 + Codex 模型）
\end{itemize}


注意：你也可以在主会话任务上设置 \texttt{model}，但这会更改共享的主会话模型。我们建议仅对隔离任务使用模型覆盖，以避免意外的上下文切换。

优先级解析顺序：

\begin{enumerate}
  \item 任务负载覆盖（最高优先级）
  \item 钩子特定默认值（例如 \texttt{hooks.gmail.model}）
  \item 智能体配置默认值
\end{enumerate}


\subsubsection{投递（渠道 + 目标）}


隔离任务可以通过顶层 \texttt{delivery} 配置投递输出：

\begin{itemize}
  \item \texttt{delivery.mode}：\texttt{announce}（投递摘要）或 \texttt{none}
  \item \texttt{delivery.channel}：\texttt{whatsapp} / \texttt{telegram} / \texttt{discord} / \texttt{slack} / \texttt{mattermost}（插件）/ \texttt{signal} / \texttt{imessage} / \texttt{last}
  \item \texttt{delivery.to}：渠道特定的接收目标
  \item \texttt{delivery.bestEffort}：投递失败时避免任务失败
\end{itemize}


当启用 announce 投递时，该轮次会抑制消息工具发送；请使用 \texttt{delivery.channel}/\texttt{delivery.to} 来指定目标。

如果省略 \texttt{delivery.channel} 或 \texttt{delivery.to}，定时任务会回退到主会话的“最后路由”（智能体最后回复的位置）。

目标格式提醒：

\begin{itemize}
  \item Slack/Discord/Mattermost（插件）目标应使用明确前缀（例如 \texttt{channel:}、\texttt{user:}）以避免歧义。
  \item Telegram 主题应使用 \texttt{:topic:} 格式（见下文）。
\end{itemize}


\paragraph{Telegram 投递目标（主题/论坛帖子）}


Telegram 通过 \texttt{message\_thread\_id} 支持论坛主题。对于定时任务投递，你可以将主题/帖子编码到 \texttt{to} 字段中：

\begin{itemize}
  \item \texttt{-1001234567890}（仅聊天 ID）
  \item \texttt{-1001234567890:topic:123}（推荐：明确的主题标记）
  \item \texttt{-1001234567890:123}（简写：数字后缀）
\end{itemize}


带前缀的目标如 \texttt{telegram:...} / \texttt{telegram:group:...} 也可接受：

\begin{itemize}
  \item \texttt{telegram:group:-1001234567890:topic:123}
\end{itemize}


\subsection{工具调用的 JSON 模式}


直接调用 Gateway网关 \texttt{cron.*} 工具（智能体工具调用或 RPC）时使用这些结构。CLI 标志接受人类可读的时间格式如 \texttt{20m}，但工具调用应使用 ISO 8601 字符串作为 \texttt{schedule.at}，并使用毫秒作为 \texttt{schedule.everyMs}。

\subsubsection{cron.add 参数}


一次性主会话任务（系统事件）：

\begin{lstlisting}[language=json]
{
  "name": "Reminder",
  "schedule": { "kind": "at", "at": "2026-02-01T16:00:00Z" },
  "sessionTarget": "main",
  "wakeMode": "now",
  "payload": { "kind": "systemEvent", "text": "Reminder text" },
  "deleteAfterRun": true
}
\end{lstlisting}


带投递的周期性隔离任务：

\begin{lstlisting}[language=json]
{
  "name": "Morning brief",
  "schedule": { "kind": "cron", "expr": "0 7 * * *", "tz": "America/Los_Angeles" },
  "sessionTarget": "isolated",
  "wakeMode": "next-heartbeat",
  "payload": {
    "kind": "agentTurn",
    "message": "Summarize overnight updates."
  },
  "delivery": {
    "mode": "announce",
    "channel": "slack",
    "to": "channel:C1234567890",
    "bestEffort": true
  }
}
\end{lstlisting}


说明：

\begin{itemize}
  \item \texttt{schedule.kind}：\texttt{at}（\texttt{at}）、\texttt{every}（\texttt{everyMs}）或 \texttt{cron}（\texttt{expr}，可选 \texttt{tz}）。
  \item \texttt{schedule.at} 接受 ISO 8601（可省略时区；省略时按 UTC 处理）。
  \item \texttt{everyMs} 为毫秒数。
  \item \texttt{sessionTarget} 必须为 \texttt{"main"} 或 \texttt{"isolated"}，且必须与 \texttt{payload.kind} 匹配。
  \item 可选字段：\texttt{agentId}、\texttt{description}、\texttt{enabled}、\texttt{deleteAfterRun}、\texttt{delivery}。
  \item \texttt{wakeMode} 省略时默认为 \texttt{"next-heartbeat"}。
\end{itemize}


\subsubsection{cron.update 参数}


\begin{lstlisting}[language=json]
{
  "jobId": "job-123",
  "patch": {
    "enabled": false,
    "schedule": { "kind": "every", "everyMs": 3600000 }
  }
}
\end{lstlisting}


说明：

\begin{itemize}
  \item \texttt{jobId} 是规范字段；\texttt{id} 可兼容使用。
  \item 在补丁中使用 \texttt{agentId: null} 可清除智能体绑定。
\end{itemize}


\subsubsection{cron.run 和 cron.remove 参数}


\begin{lstlisting}[language=json]
{ "jobId": "job-123", "mode": "force" }
\end{lstlisting}


\begin{lstlisting}[language=json]
{ "jobId": "job-123" }
\end{lstlisting}


\subsection{存储与历史}


\begin{itemize}
  \item 任务存储：\texttt{\textasciitilde{}/.openclaw/cron/jobs.json}（Gateway网关管理的 JSON）。
  \item 运行历史：\texttt{\textasciitilde{}/.openclaw/cron/runs/.jsonl}（JSONL，自动清理）。
  \item 覆盖存储路径：配置中的 \texttt{cron.store}。
\end{itemize}


\subsection{配置}


\begin{lstlisting}[language=json]
{
  cron: {
    enabled: true, // 默认 true
    store: "~/.openclaw/cron/jobs.json",
    maxConcurrentRuns: 1, // 默认 1
  },
}
\end{lstlisting}


完全禁用定时任务：

\begin{itemize}
  \item \texttt{cron.enabled: false}（配置）
  \item \texttt{OPENCLAW\_SKIP\_CRON=1}（环境变量）
\end{itemize}


\subsection{CLI 快速开始}


一次性提醒（UTC ISO，成功后自动删除）：

\begin{lstlisting}[language=bash]
openclaw cron add \
  --name "Send reminder" \
  --at "2026-01-12T18:00:00Z" \
  --session main \
  --system-event "Reminder: submit expense report." \
  --wake now \
  --delete-after-run
\end{lstlisting}


一次性提醒（主会话，立即唤醒）：

\begin{lstlisting}[language=bash]
openclaw cron add \
  --name "Calendar check" \
  --at "20m" \
  --session main \
  --system-event "Next heartbeat: check calendar." \
  --wake now
\end{lstlisting}


周期性隔离任务（投递到 WhatsApp）：

\begin{lstlisting}[language=bash]
openclaw cron add \
  --name "Morning status" \
  --cron "0 7 * * *" \
  --tz "America/Los_Angeles" \
  --session isolated \
  --message "Summarize inbox + calendar for today." \
  --announce \
  --channel whatsapp \
  --to "+15551234567"
\end{lstlisting}


周期性隔离任务（投递到 Telegram 主题）：

\begin{lstlisting}[language=bash]
openclaw cron add \
  --name "Nightly summary (topic)" \
  --cron "0 22 * * *" \
  --tz "America/Los_Angeles" \
  --session isolated \
  --message "Summarize today; send to the nightly topic." \
  --announce \
  --channel telegram \
  --to "-1001234567890:topic:123"
\end{lstlisting}


带模型和思维覆盖的隔离任务：

\begin{lstlisting}[language=bash]
openclaw cron add \
  --name "Deep analysis" \
  --cron "0 6 * * 1" \
  --tz "America/Los_Angeles" \
  --session isolated \
  --message "Weekly deep analysis of project progress." \
  --model "opus" \
  --thinking high \
  --announce \
  --channel whatsapp \
  --to "+15551234567"
\end{lstlisting}


智能体选择（多智能体配置）：

\begin{lstlisting}[language=bash]
# 将任务绑定到智能体 "ops"（如果该智能体不存在则回退到默认智能体）
openclaw cron add --name "Ops sweep" --cron "0 6 * * *" --session isolated --message "Check ops queue" --agent ops

# 切换或清除现有任务的智能体
openclaw cron edit <jobId> --agent ops
openclaw cron edit <jobId> --clear-agent
\end{lstlisting}


手动运行（调试）：

\begin{lstlisting}[language=bash]
openclaw cron run <jobId> --force
\end{lstlisting}


编辑现有任务（补丁字段）：

\begin{lstlisting}[language=bash]
openclaw cron edit <jobId> \
  --message "Updated prompt" \
  --model "opus" \
  --thinking low
\end{lstlisting}


运行历史：

\begin{lstlisting}[language=bash]
openclaw cron runs --id <jobId> --limit 50
\end{lstlisting}


不创建任务直接发送系统事件：

\begin{lstlisting}[language=bash]
openclaw system event --mode now --text "Next heartbeat: check battery."
\end{lstlisting}


\subsection{Gateway网关 API 接口}


\begin{itemize}
  \item \texttt{cron.list}、\texttt{cron.status}、\texttt{cron.add}、\texttt{cron.update}、\texttt{cron.remove}
  \item \texttt{cron.run}（强制或到期）、\texttt{cron.runs}
\end{itemize}

如需不创建任务直接发送系统事件，请使用 \href{/cli/system}{\texttt{openclaw system event}}。

\subsection{故障排除}


\subsubsection{"没有任何任务运行"}


\begin{itemize}
  \item 检查定时任务是否已启用：\texttt{cron.enabled} 和 \texttt{OPENCLAW\_SKIP\_CRON}。
  \item 检查 Gateway网关是否持续运行（定时任务运行在 Gateway网关进程内部）。
  \item 对于 \texttt{cron} 调度：确认时区（\texttt{--tz}）与主机时区的关系。
\end{itemize}


\subsubsection{Telegram 投递到了错误的位置}


\begin{itemize}
  \item 对于论坛主题，使用 \texttt{-100…:topic:} 以确保明确无歧义。
  \item 如果你在日志或存储的"最后路由"目标中看到 \texttt{telegram:...} 前缀，这是正常的；定时任务投递接受这些前缀并仍能正确解析主题 ID。
\end{itemize}



\clearpage

% === 源文件: automation/cron-vs-heartbeat.md ===
\section{定时任务与心跳：何时使用哪种方式}


心跳和定时任务都可以按计划运行任务。本指南帮助你根据使用场景选择合适的机制。

\subsection{快速决策指南}



\begin{table}[h]
\centering
\small
\begin{tabular}{|l|l|l|}
\hline
使用场景 & 推荐方式 & 原因 \\
\hline
每 30 分钟检查收件箱 & 心跳 & 可与其他检查批量处理，具备上下文感知能力 \\
\hline
每天上午 9 点准时发送报告 & 定时任务（隔离式） & 需要精确定时 \\
\hline
监控日历中即将到来的事件 & 心跳 & 天然适合周期性感知 \\
\hline
运行每周深度分析 & 定时任务（隔离式） & 独立任务，可使用不同模型 \\
\hline
20 分钟后提醒我 & 定时任务（主会话，\texttt{--at}） & 精确定时的一次性任务 \\
\hline
后台项目健康检查 & 心跳 & 搭载在现有周期上 \\
\hline
\end{tabular}
\end{table}



\subsection{心跳：周期性感知}


心跳在\textbf{主会话}中以固定间隔运行（默认：30 分钟）。它的设计目的是让智能体检查各种事项并呈现重要信息。

\subsubsection{何时使用心跳}


\begin{itemize}
  \item \textbf{多个周期性检查}：与其设置 5 个独立的定时任务分别检查收件箱、日历、天气、通知和项目状态，不如用一次心跳批量处理所有内容。
  \item \textbf{上下文感知决策}：智能体拥有完整的主会话上下文，因此可以智能判断哪些紧急、哪些可以等待。
  \item \textbf{对话连续性}：心跳运行共享同一会话，因此智能体记得最近的对话，可以自然地进行后续跟进。
  \item \textbf{低开销监控}：一次心跳替代多个小型轮询任务。
\end{itemize}


\subsubsection{心跳优势}


\begin{itemize}
  \item \textbf{批量处理多项检查}：一次智能体轮次可以同时审查收件箱、日历和通知。
  \item \textbf{减少 API 调用}：一次心跳比 5 个隔离式定时任务更经济。
  \item \textbf{上下文感知}：智能体了解你一直在做什么，可以据此排定优先级。
  \item \textbf{智能抑制}：如果没有需要关注的事项，智能体回复 \texttt{HEARTBEAT\_OK}，不会投递任何消息。
  \item \textbf{自然定时}：会根据队列负载略有漂移，但对大多数监控来说没有问题。
\end{itemize}


\subsubsection{心跳示例：HEARTBEAT.md 检查清单}


\begin{lstlisting}
# Heartbeat checklist

- Check email for urgent messages
- Review calendar for events in next 2 hours
- If a background task finished, summarize results
- If idle for 8+ hours, send a brief check-in
\end{lstlisting}


智能体在每次心跳时读取此清单，并在一次轮次中处理所有项目。

\subsubsection{配置心跳}


\begin{lstlisting}[language=json]
{
  agents: {
    defaults: {
      heartbeat: {
        every: "30m", // 间隔
        target: "last", // 告警投递目标
        activeHours: { start: "08:00", end: "22:00" }, // 可选
      },
    },
  },
}
\end{lstlisting}


完整配置请参阅\href{/gateway/heartbeat}{心跳}。

\subsection{定时任务：精确调度}


定时任务在\textbf{精确时间}运行，可以在隔离会话中运行而不影响主会话上下文。

\subsubsection{何时使用定时任务}


\begin{itemize}
  \item \textbf{需要精确定时}："每周一上午 9:00 发送"（而不是"大约 9 点左右"）。
  \item \textbf{独立任务}：不需要对话上下文的任务。
  \item \textbf{不同的模型/思维级别}：需要更强大模型的深度分析。
  \item \textbf{一次性提醒}：使用 \texttt{--at} 实现"20 分钟后提醒我"。
  \item \textbf{嘈杂/频繁的任务}：会把主会话历史搞得杂乱的任务。
  \item \textbf{外部触发器}：无论智能体是否处于活跃状态都应独立运行的任务。
\end{itemize}


\subsubsection{定时任务优势}


\begin{itemize}
  \item \textbf{精确定时}：支持带时区的 5 字段 cron 表达式。
  \item \textbf{会话隔离}：在 \texttt{cron:} 中运行，不会污染主会话历史。
  \item \textbf{模型覆盖}：可按任务使用更便宜或更强大的模型。
  \item \textbf{投递控制}：隔离任务默认以 \texttt{announce} 投递摘要，可选 \texttt{none} 仅内部运行。
  \item \textbf{无需智能体上下文}：即使主会话空闲或已压缩，也能运行。
  \item \textbf{一次性支持}：\texttt{--at} 用于精确的未来时间戳。
\end{itemize}


\subsubsection{定时任务示例：每日早间简报}


\begin{lstlisting}[language=bash]
openclaw cron add \
  --name "Morning briefing" \
  --cron "0 7 * * *" \
  --tz "America/New_York" \
  --session isolated \
  --message "Generate today's briefing: weather, calendar, top emails, news summary." \
  --model opus \
  --announce \
  --channel whatsapp \
  --to "+15551234567"
\end{lstlisting}


这会在纽约时间每天早上 7:00 准时运行，使用 Opus 保证质量，并直接投递到 WhatsApp。

\subsubsection{定时任务示例：一次性提醒}


\begin{lstlisting}[language=bash]
openclaw cron add \
  --name "Meeting reminder" \
  --at "20m" \
  --session main \
  --system-event "Reminder: standup meeting starts in 10 minutes." \
  --wake now \
  --delete-after-run
\end{lstlisting}


完整 CLI 参考请参阅\href{/automation/cron-jobs}{定时任务}。

\subsection{决策流程图}


\begin{lstlisting}
任务是否需要在精确时间运行？
  是 -> 使用定时任务
  否 -> 继续...

任务是否需要与主会话隔离？
  是 -> 使用定时任务（隔离式）
  否 -> 继续...

此任务能否与其他周期性检查批量处理？
  是 -> 使用心跳（添加到 HEARTBEAT.md）
  否 -> 使用定时任务

这是一次性提醒吗？
  是 -> 使用定时任务配合 --at
  否 -> 继续...

是否需要不同的模型或思维级别？
  是 -> 使用定时任务（隔离式）配合 --model/--thinking
  否 -> 使用心跳
\end{lstlisting}


\subsection{组合使用}


最高效的配置是\textbf{两者结合}：

\begin{enumerate}
  \item \textbf{心跳}处理常规监控（收件箱、日历、通知），每 30 分钟批量处理一次。
  \item \textbf{定时任务}处理精确调度（每日报告、每周回顾）和一次性提醒。
\end{enumerate}


\subsubsection{示例：高效自动化配置}


\textbf{HEARTBEAT.md}（每 30 分钟检查一次）：

\begin{lstlisting}
# Heartbeat checklist

- Scan inbox for urgent emails
- Check calendar for events in next 2h
- Review any pending tasks
- Light check-in if quiet for 8+ hours
\end{lstlisting}


\textbf{定时任务}（精确定时）：

\begin{lstlisting}[language=bash]
# 每天早上 7 点的早间简报
openclaw cron add --name "Morning brief" --cron "0 7 * * *" --session isolated --message "..." --announce

# 每周一上午 9 点的项目回顾
openclaw cron add --name "Weekly review" --cron "0 9 * * 1" --session isolated --message "..." --model opus

# 一次性提醒
openclaw cron add --name "Call back" --at "2h" --session main --system-event "Call back the client" --wake now
\end{lstlisting}


\subsection{Lobster：带审批的确定性工作流}


Lobster 是用于\textbf{多步骤工具管道}的工作流运行时，适用于需要确定性执行和明确审批的场景。当任务不只是单次智能体轮次，且你需要可恢复的带人工检查点的工作流时，使用它。

\subsubsection{何时适合使用 Lobster}


\begin{itemize}
  \item \textbf{多步骤自动化}：你需要一个固定的工具调用管道，而不是一次性提示。
  \item \textbf{审批关卡}：副作用应暂停直到你批准，然后继续执行。
  \item \textbf{可恢复运行}：继续暂停的工作流而无需重新运行之前的步骤。
\end{itemize}


\subsubsection{如何与心跳和定时任务配合}


\begin{itemize}
  \item \textbf{心跳/定时任务}决定\textit{何时}运行。
  \item \textbf{Lobster} 定义运行开始后\textit{执行哪些步骤}。
\end{itemize}


对于计划性工作流，使用定时任务或心跳触发一次调用 Lobster 的智能体轮次。对于临时工作流，直接调用 Lobster。

\subsubsection{操作说明（来自代码）}


\begin{itemize}
  \item Lobster 以\textbf{本地子进程}（\texttt{lobster} CLI）在工具模式下运行，并返回 \textbf{JSON 信封}。
  \item 如果工具返回 \texttt{needs\_approval}，你需要使用 \texttt{resumeToken} 和 \texttt{approve} 标志来恢复。
  \item 该工具是\textbf{可选插件}；建议通过 \texttt{tools.alsoAllow: ["lobster"]} 附加启用。
  \item 如果传入 \texttt{lobsterPath}，必须是\textbf{绝对路径}。
\end{itemize}


完整用法和示例请参阅 \href{/tools/lobster}{Lobster}。

\subsection{主会话与隔离会话}


心跳和定时任务都可以与主会话交互，但方式不同：


\begin{table}[h]
\centering
\small
\begin{tabular}{|l|l|l|l|}
\hline
 & 心跳 & 定时任务（主会话） & 定时任务（隔离式） \\
\hline
会话 & 主会话 & 主会话（通过系统事件） & \texttt{cron:} \\
\hline
历史 & 共享 & 共享 & 每次运行全新 \\
\hline
上下文 & 完整 & 完整 & 无（从零开始） \\
\hline
模型 & 主会话模型 & 主会话模型 & 可覆盖 \\
\hline
输出 & 非 \texttt{HEARTBEAT\_OK} 时投递 & 心跳提示 + 事件 & announce 摘要（默认） \\
\hline
\end{tabular}
\end{table}



\subsubsection{何时使用主会话定时任务}


当你需要以下场景时，使用 \texttt{--session main} 配合 \texttt{--system-event}：

\begin{itemize}
  \item 提醒/事件出现在主会话上下文中
  \item 智能体在下一次心跳时带着完整上下文处理它
  \item 不需要单独的隔离运行
\end{itemize}


\begin{lstlisting}[language=bash]
openclaw cron add \
  --name "Check project" \
  --every "4h" \
  --session main \
  --system-event "Time for a project health check" \
  --wake now
\end{lstlisting}


\subsubsection{何时使用隔离式定时任务}


当你需要以下场景时，使用 \texttt{--session isolated}：

\begin{itemize}
  \item 无先前上下文的全新环境
  \item 不同的模型或思维设置
  \item 输出可通过 \texttt{announce} 直接投递摘要（或用 \texttt{none} 仅内部运行）
  \item 不会把主会话搞得杂乱的历史记录
\end{itemize}


\begin{lstlisting}[language=bash]
openclaw cron add \
  --name "Deep analysis" \
  --cron "0 6 * * 0" \
  --session isolated \
  --message "Weekly codebase analysis..." \
  --model opus \
  --thinking high \
  --announce
\end{lstlisting}


\subsection{成本考量}



\begin{table}[h]
\centering
\small
\begin{tabular}{|l|l|}
\hline
机制 & 成本特征 \\
\hline
心跳 & 每 N 分钟一次轮次；随 HEARTBEAT.md 大小扩展 \\
\hline
定时任务（主会话） & 将事件添加到下一次心跳（无隔离轮次） \\
\hline
定时任务（隔离式） & 每个任务一次完整智能体轮次；可使用更便宜的模型 \\
\hline
\end{tabular}
\end{table}



\textbf{建议}：

\begin{itemize}
  \item 保持 \texttt{HEARTBEAT.md} 精简以减少 token 开销。
  \item 将类似的检查批量放入心跳，而不是创建多个定时任务。
  \item 如果只需要内部处理，在心跳上使用 \texttt{target: "none"}。
  \item 对常规任务使用隔离式定时任务配合更便宜的模型。
\end{itemize}


\subsection{相关内容}


\begin{itemize}
  \item \href{/gateway/heartbeat}{心跳} - 完整的心跳配置
  \item \href{/automation/cron-jobs}{定时任务} - 完整的定时任务 CLI 和 API 参考
  \item \href{/cli/system}{系统} - 系统事件 + 心跳控制
\end{itemize}



\clearpage

% === 源文件: automation/gmail-pubsub.md ===
\section{Gmail Pub/Sub -> OpenClaw}


目标：Gmail watch -> Pub/Sub 推送 -> \texttt{gog gmail watch serve} -> OpenClaw webhook。

\subsection{前置条件}


\begin{itemize}
  \item 已安装并登录 \texttt{gcloud}（\href{https://docs.cloud.google.com/sdk/docs/install-sdk}{安装指南}）。
  \item 已安装 \texttt{gog} (gogcli) 并为 Gmail 账户授权（\href{https://gogcli.sh/}{gogcli.sh}）。
  \item 已启用 OpenClaw hooks（参见 \href{/automation/webhook}{Webhooks}）。
  \item 已登录 \texttt{tailscale}（\href{https://tailscale.com/}{tailscale.com}）。支持的设置使用 Tailscale Funnel 作为公共 HTTPS 端点。
\end{itemize}

其他隧道服务也可以使用，但需要自行配置/不受支持，需要手动接入。
目前，我们支持的是 Tailscale。

示例 hook 配置（启用 Gmail 预设映射）：

\begin{lstlisting}[language=json]
{
  hooks: {
    enabled: true,
    token: "OPENCLAW_HOOK_TOKEN",
    path: "/hooks",
    presets: ["gmail"],
  },
}
\end{lstlisting}


要将 Gmail 摘要投递到聊天界面，请用设置了 \texttt{deliver} 以及可选的 \texttt{channel}/\texttt{to} 的映射覆盖预设：

\begin{lstlisting}[language=json]
{
  hooks: {
    enabled: true,
    token: "OPENCLAW_HOOK_TOKEN",
    presets: ["gmail"],
    mappings: [
      {
        match: { path: "gmail" },
        action: "agent",
        wakeMode: "now",
        name: "Gmail",
        sessionKey: "hook:gmail:{{messages[0].id}}",
        messageTemplate: "New email from {{messages[0].from}}\nSubject: {{messages[0].subject}}\n{{messages[0].snippet}}\n{{messages[0].body}}",
        model: "openai/gpt-5.2-mini",
        deliver: true,
        channel: "last",
        // to: "+15551234567"
      },
    ],
  },
}
\end{lstlisting}


如果你想使用固定渠道，请设置 \texttt{channel} + \texttt{to}。否则 \texttt{channel: "last"} 会使用上次的投递路由（默认回退到 WhatsApp）。

要为 Gmail 运行强制使用更便宜的模型，请在映射中设置 \texttt{model}（\texttt{provider/model} 或别名）。如果你强制启用了 \texttt{agents.defaults.models}，请将其包含在内。

要专门为 Gmail hooks 设置默认模型和思考级别，请在配置中添加 \texttt{hooks.gmail.model} / \texttt{hooks.gmail.thinking}：

\begin{lstlisting}[language=json]
{
  hooks: {
    gmail: {
      model: "openrouter/meta-llama/llama-3.3-70b-instruct:free",
      thinking: "off",
    },
  },
}
\end{lstlisting}


注意事项：

\begin{itemize}
  \item 映射中的每个 hook 的 \texttt{model}/\texttt{thinking} 仍会覆盖这些默认值。
  \item 回退顺序：\texttt{hooks.gmail.model} → \texttt{agents.defaults.model.fallbacks} → 主模型（认证/速率限制/超时）。
  \item 如果设置了 \texttt{agents.defaults.models}，Gmail 模型必须在允许列表中。
  \item Gmail hook 内容默认使用外部内容安全边界包装。
\end{itemize}

要禁用（危险），请设置 \texttt{hooks.gmail.allowUnsafeExternalContent: true}。

要进一步自定义负载处理，请添加 \texttt{hooks.mappings} 或在 \texttt{hooks.transformsDir} 下添加 JS/TS 转换模块（参见 \href{/automation/webhook}{Webhooks}）。

\subsection{向导（推荐）}


使用 OpenClaw 助手将所有内容接入在一起（在 macOS 上通过 brew 安装依赖）：

\begin{lstlisting}[language=bash]
openclaw webhooks gmail setup \
  --account openclaw@gmail.com
\end{lstlisting}


默认设置：

\begin{itemize}
  \item 使用 Tailscale Funnel 作为公共推送端点。
  \item 为 \texttt{openclaw webhooks gmail run} 写入 \texttt{hooks.gmail} 配置。
  \item 启用 Gmail hook 预设（\texttt{hooks.presets: ["gmail"]}）。
\end{itemize}


路径说明：当启用 \texttt{tailscale.mode} 时，OpenClaw 会自动将 \texttt{hooks.gmail.serve.path} 设置为 \texttt{/}，并将公共路径保持在 \texttt{hooks.gmail.tailscale.path}（默认 \texttt{/gmail-pubsub}），因为 Tailscale 在代理之前会剥离设置的路径前缀。
如果你需要后端接收带前缀的路径，请将 \texttt{hooks.gmail.tailscale.target}（或 \texttt{--tailscale-target}）设置为完整 URL，如 \texttt{http://127.0.0.1:8788/gmail-pubsub}，并匹配 \texttt{hooks.gmail.serve.path}。

想要自定义端点？使用 \texttt{--push-endpoint } 或 \texttt{--tailscale off}。

平台说明：在 macOS 上，向导通过 Homebrew 安装 \texttt{gcloud}、\texttt{gogcli} 和 \texttt{tailscale}；在 Linux 上请先手动安装它们。

Gateway 网关自动启动（推荐）：

\begin{itemize}
  \item 当 \texttt{hooks.enabled=true} 且设置了 \texttt{hooks.gmail.account} 时，Gateway 网关会在启动时运行 \texttt{gog gmail watch serve} 并自动续期 watch。
  \item 设置 \texttt{OPENCLAW\_SKIP\_GMAIL\_WATCHER=1} 可退出（如果你自己运行守护进程则很有用）。
  \item 不要同时运行手动守护进程，否则会遇到 \texttt{listen tcp 127.0.0.1:8788: bind: address already in use}。
\end{itemize}


手动守护进程（启动 \texttt{gog gmail watch serve} + 自动续期）：

\begin{lstlisting}[language=bash]
openclaw webhooks gmail run
\end{lstlisting}


\subsection{一次性设置}


\begin{enumerate}
  \item 选择\textbf{拥有 \texttt{gog} 使用的 OAuth 客户端}的 GCP 项目。
\end{enumerate}


\begin{lstlisting}[language=bash]
gcloud auth login
gcloud config set project <project-id>
\end{lstlisting}


注意：Gmail watch 要求 Pub/Sub 主题与 OAuth 客户端位于同一项目中。

\begin{enumerate}
  \item 启用 API：
\end{enumerate}


\begin{lstlisting}[language=bash]
gcloud services enable gmail.googleapis.com pubsub.googleapis.com
\end{lstlisting}


\begin{enumerate}
  \item 创建主题：
\end{enumerate}


\begin{lstlisting}[language=bash]
gcloud pubsub topics create gog-gmail-watch
\end{lstlisting}


\begin{enumerate}
  \item 允许 Gmail push 发布：
\end{enumerate}


\begin{lstlisting}[language=bash]
gcloud pubsub topics add-iam-policy-binding gog-gmail-watch \
  --member=serviceAccount:gmail-api-push@system.gserviceaccount.com \
  --role=roles/pubsub.publisher
\end{lstlisting}


\subsection{启动 watch}


\begin{lstlisting}[language=bash]
gog gmail watch start \
  --account openclaw@gmail.com \
  --label INBOX \
  --topic projects/<project-id>/topics/gog-gmail-watch
\end{lstlisting}


保存输出中的 \texttt{history\_id}（用于调试）。

\subsection{运行推送处理程序}


本地示例（共享 token 认证）：

\begin{lstlisting}[language=bash]
gog gmail watch serve \
  --account openclaw@gmail.com \
  --bind 127.0.0.1 \
  --port 8788 \
  --path /gmail-pubsub \
  --token <shared> \
  --hook-url http://127.0.0.1:18789/hooks/gmail \
  --hook-token OPENCLAW_HOOK_TOKEN \
  --include-body \
  --max-bytes 20000
\end{lstlisting}


注意事项：

\begin{itemize}
  \item \texttt{--token} 保护推送端点（\texttt{x-gog-token} 或 \texttt{?token=}）。
  \item \texttt{--hook-url} 指向 OpenClaw \texttt{/hooks/gmail}（已映射；隔离运行 + 摘要发送到主线程）。
  \item \texttt{--include-body} 和 \texttt{--max-bytes} 控制发送到 OpenClaw 的正文片段。
\end{itemize}


推荐：\texttt{openclaw webhooks gmail run} 封装了相同的流程并自动续期 watch。

\subsection{暴露处理程序（高级，不受支持）}


如果你需要非 Tailscale 隧道，请手动接入并在推送订阅中使用公共 URL（不受支持，无保护措施）：

\begin{lstlisting}[language=bash]
cloudflared tunnel --url http://127.0.0.1:8788 --no-autoupdate
\end{lstlisting}


使用生成的 URL 作为推送端点：

\begin{lstlisting}[language=bash]
gcloud pubsub subscriptions create gog-gmail-watch-push \
  --topic gog-gmail-watch \
  --push-endpoint "https://<public-url>/gmail-pubsub?token=<shared>"
\end{lstlisting}


生产环境：使用稳定的 HTTPS 端点并配置 Pub/Sub OIDC JWT，然后运行：

\begin{lstlisting}[language=bash]
gog gmail watch serve --verify-oidc --oidc-email <svc@...>
\end{lstlisting}


\subsection{测试}


向被监视的收件箱发送一条消息：

\begin{lstlisting}[language=bash]
gog gmail send \
  --account openclaw@gmail.com \
  --to openclaw@gmail.com \
  --subject "watch test" \
  --body "ping"
\end{lstlisting}


检查 watch 状态和历史记录：

\begin{lstlisting}[language=bash]
gog gmail watch status --account openclaw@gmail.com
gog gmail history --account openclaw@gmail.com --since <historyId>
\end{lstlisting}


\subsection{故障排除}


\begin{itemize}
  \item \texttt{Invalid topicName}：项目不匹配（主题不在 OAuth 客户端项目中）。
  \item \texttt{User not authorized}：主题缺少 \texttt{roles/pubsub.publisher}。
  \item 空消息：Gmail push 仅提供 \texttt{historyId}；通过 \texttt{gog gmail history} 获取。
\end{itemize}


\subsection{清理}


\begin{lstlisting}[language=bash]
gog gmail watch stop --account openclaw@gmail.com
gcloud pubsub subscriptions delete gog-gmail-watch-push
gcloud pubsub topics delete gog-gmail-watch
\end{lstlisting}



\clearpage

% === 源文件: automation/hooks.md ===
\section{Hooks}


Hooks 提供了一个可扩展的事件驱动系统，用于响应智能体命令和事件自动执行操作。Hooks 从目录中自动发现，可以通过 CLI 命令管理，类似于 OpenClaw 中 Skills 的工作方式。

\subsection{入门指南}


Hooks 是在事件发生时运行的小脚本。有两种类型：

\begin{itemize}
  \item \textbf{Hooks}（本页）：当智能体事件触发时在 Gateway 网关内运行，如 \texttt{/new}、\texttt{/reset}、\texttt{/stop} 或生命周期事件。
  \item \textbf{Webhooks}：外部 HTTP webhooks，让其他系统触发 OpenClaw 中的工作。参见 \href{/automation/webhook}{Webhook Hooks} 或使用 \texttt{openclaw webhooks} 获取 Gmail 助手命令。
\end{itemize}


Hooks 也可以捆绑在插件中；参见 \href{/tools/plugin\#plugin-hooks}{插件}。

常见用途：

\begin{itemize}
  \item 重置会话时保存记忆快照
  \item 保留命令审计跟踪用于故障排除或合规
  \item 会话开始或结束时触发后续自动化
  \item 事件触发时向智能体工作区写入文件或调用外部 API
\end{itemize}


如果你能写一个小的 TypeScript 函数，你就能写一个 hook。Hooks 会自动发现，你可以通过 CLI 启用或禁用它们。

\subsection{概述}


hooks 系统允许你：

\begin{itemize}
  \item 在发出 \texttt{/new} 时将会话上下文保存到记忆
  \item 记录所有命令以供审计
  \item 在智能体生命周期事件上触发自定义自动化
  \item 在不修改核心代码的情况下扩展 OpenClaw 的行为
\end{itemize}


\subsection{入门}


\subsubsection{捆绑的 Hooks}


OpenClaw 附带三个自动发现的捆绑 hooks：

\begin{itemize}
  \item \textbf{💾 session-memory}：当你发出 \texttt{/new} 时将会话上下文保存到智能体工作区（默认 \texttt{\textasciitilde{}/.openclaw/workspace/memory/}）
  \item \textbf{📝 command-logger}：将所有命令事件记录到 \texttt{\textasciitilde{}/.openclaw/logs/commands.log}
  \item \textbf{🚀 boot-md}：当 Gateway 网关启动时运行 \texttt{BOOT.md}（需要启用内部 hooks）
\end{itemize}


列出可用的 hooks：

\begin{lstlisting}[language=bash]
openclaw hooks list
\end{lstlisting}


启用一个 hook：

\begin{lstlisting}[language=bash]
openclaw hooks enable session-memory
\end{lstlisting}


检查 hook 状态：

\begin{lstlisting}[language=bash]
openclaw hooks check
\end{lstlisting}


获取详细信息：

\begin{lstlisting}[language=bash]
openclaw hooks info session-memory
\end{lstlisting}


\subsubsection{新手引导}


在新手引导期间（\texttt{openclaw onboard}），你将被提示启用推荐的 hooks。向导会自动发现符合条件的 hooks 并呈现供选择。

\subsection{Hook 发现}


Hooks 从三个目录自动发现（按优先级顺序）：

\begin{enumerate}
  \item \textbf{工作区 hooks}：\texttt{/hooks/}（每智能体，最高优先级）
  \item \textbf{托管 hooks}：\texttt{\textasciitilde{}/.openclaw/hooks/}（用户安装，跨工作区共享）
  \item \textbf{捆绑 hooks}：\texttt{/dist/hooks/bundled/}（随 OpenClaw 附带）
\end{enumerate}


托管 hook 目录可以是\textbf{单个 hook} 或 \textbf{hook 包}（包目录）。

每个 hook 是一个包含以下内容的目录：

\begin{lstlisting}
my-hook/
├── HOOK.md          # 元数据 + 文档
└── handler.ts       # 处理程序实现
\end{lstlisting}


\subsection{Hook 包（npm/archives）}


Hook 包是标准的 npm 包，通过 \texttt{package.json} 中的 \texttt{openclaw.hooks} 导出一个或多个 hooks。使用以下命令安装：

\begin{lstlisting}[language=bash]
openclaw hooks install <path-or-spec>
\end{lstlisting}


示例 \texttt{package.json}：

\begin{lstlisting}[language=json]
{
  "name": "@acme/my-hooks",
  "version": "0.1.0",
  "openclaw": {
    "hooks": ["./hooks/my-hook", "./hooks/other-hook"]
  }
}
\end{lstlisting}


每个条目指向包含 \texttt{HOOK.md} 和 \texttt{handler.ts}（或 \texttt{index.ts}）的 hook 目录。
Hook 包可以附带依赖；它们将安装在 \texttt{\textasciitilde{}/.openclaw/hooks/} 下。

\subsection{Hook 结构}


\subsubsection{HOOK.md 格式}


\texttt{HOOK.md} 文件在 YAML frontmatter 中包含元数据，加上 Markdown 文档：

\begin{lstlisting}
---
name: my-hook
description: "Short description of what this hook does"
homepage: https://docs.openclaw.ai/automation/hooks#my-hook
metadata:
  { "openclaw": { "emoji": "🔗", "events": ["command:new"], "requires": { "bins": ["node"] } } }
---

# My Hook

Detailed documentation goes here...

## What It Does

- Listens for `/new` commands
- Performs some action
- Logs the result

## Requirements

- Node.js must be installed

## Configuration

No configuration needed.
\end{lstlisting}


\subsubsection{元数据字段}


\texttt{metadata.openclaw} 对象支持：

\begin{itemize}
  \item \textbf{\texttt{emoji}}：CLI 的显示表情符号（例如 \texttt{"💾"}）
  \item \textbf{\texttt{events}}：要监听的事件数组（例如 \texttt{["command:new", "command:reset"]}）
  \item \textbf{\texttt{export}}：要使用的命名导出（默认为 \texttt{"default"}）
  \item \textbf{\texttt{homepage}}：文档 URL
  \item \textbf{\texttt{requires}}：可选要求
  \item \textbf{\texttt{bins}}：PATH 中需要的二进制文件（例如 \texttt{["git", "node"]}）
  \item \textbf{\texttt{anyBins}}：这些二进制文件中至少有一个必须存在
  \item \textbf{\texttt{env}}：需要的环境变量
  \item \textbf{\texttt{config}}：需要的配置路径（例如 \texttt{["workspace.dir"]}）
  \item \textbf{\texttt{os}}：需要的平台（例如 \texttt{["darwin", "linux"]}）
  \item \textbf{\texttt{always}}：绕过资格检查（布尔值）
  \item \textbf{\texttt{install}}：安装方法（对于捆绑 hooks：\texttt{[\{"id":"bundled","kind":"bundled"\}]}）
\end{itemize}


\subsubsection{处理程序实现}


\texttt{handler.ts} 文件导出一个 \texttt{HookHandler} 函数：

\begin{lstlisting}[language=typescript]
import type { HookHandler } from "../../src/hooks/hooks.js";

const myHandler: HookHandler = async (event) => {
  // Only trigger on 'new' command
  if (event.type !== "command" || event.action !== "new") {
    return;
  }

  console.log(`[my-hook] New command triggered`);
  console.log(`  Session: ${event.sessionKey}`);
  console.log(`  Timestamp: ${event.timestamp.toISOString()}`);

  // Your custom logic here

  // Optionally send message to user
  event.messages.push("✨ My hook executed!");
};

export default myHandler;
\end{lstlisting}


\paragraph{事件上下文}


每个事件包含：

\begin{lstlisting}[language=typescript]
{
  type: 'command' | 'session' | 'agent' | 'gateway',
  action: string,              // e.g., 'new', 'reset', 'stop'
  sessionKey: string,          // Session identifier
  timestamp: Date,             // When the event occurred
  messages: string[],          // Push messages here to send to user
  context: {
    sessionEntry?: SessionEntry,
    sessionId?: string,
    sessionFile?: string,
    commandSource?: string,    // e.g., 'whatsapp', 'telegram'
    senderId?: string,
    workspaceDir?: string,
    bootstrapFiles?: WorkspaceBootstrapFile[],
    cfg?: OpenClawConfig
  }
}
\end{lstlisting}


\subsection{事件类型}


\subsubsection{命令事件}


当发出智能体命令时触发：

\begin{itemize}
  \item \textbf{\texttt{command}}：所有命令事件（通用监听器）
  \item \textbf{\texttt{command:new}}：当发出 \texttt{/new} 命令时
  \item \textbf{\texttt{command:reset}}：当发出 \texttt{/reset} 命令时
  \item \textbf{\texttt{command:stop}}：当发出 \texttt{/stop} 命令时
\end{itemize}


\subsubsection{智能体事件}


\begin{itemize}
  \item \textbf{\texttt{agent:bootstrap}}：在注入工作区引导文件之前（hooks 可以修改 \texttt{context.bootstrapFiles}）
\end{itemize}


\subsubsection{Gateway 网关事件}


当 Gateway 网关启动时触发：

\begin{itemize}
  \item \textbf{\texttt{gateway:startup}}：在渠道启动和 hooks 加载之后
\end{itemize}


\subsubsection{工具结果 Hooks（插件 API）}


这些 hooks 不是事件流监听器；它们让插件在 OpenClaw 持久化工具结果之前同步调整它们。

\begin{itemize}
  \item \textbf{\texttt{tool\_result\_persist}}：在工具结果写入会话记录之前转换它们。必须是同步的；返回更新后的工具结果负载或 \texttt{undefined} 保持原样。参见 \href{/concepts/agent-loop}{智能体循环}。
\end{itemize}


\subsubsection{未来事件}


计划中的事件类型：

\begin{itemize}
  \item \textbf{\texttt{session:start}}：当新会话开始时
  \item \textbf{\texttt{session:end}}：当会话结束时
  \item \textbf{\texttt{agent:error}}：当智能体遇到错误时
  \item \textbf{\texttt{message:sent}}：当消息被发送时
  \item \textbf{\texttt{message:received}}：当消息被接收时
\end{itemize}


\subsection{创建自定义 Hooks}


\subsubsection{1. 选择位置}


\begin{itemize}
  \item \textbf{工作区 hooks}（\texttt{/hooks/}）：每智能体，最高优先级
  \item \textbf{托管 hooks}（\texttt{\textasciitilde{}/.openclaw/hooks/}）：跨工作区共享
\end{itemize}


\subsubsection{2. 创建目录结构}


\begin{lstlisting}[language=bash]
mkdir -p ~/.openclaw/hooks/my-hook
cd ~/.openclaw/hooks/my-hook
\end{lstlisting}


\subsubsection{3. 创建 HOOK.md}


\begin{lstlisting}
---
name: my-hook
description: "Does something useful"
metadata: { "openclaw": { "emoji": "🎯", "events": ["command:new"] } }
---

# My Custom Hook

This hook does something useful when you issue `/new`.
\end{lstlisting}


\subsubsection{4. 创建 handler.ts}


\begin{lstlisting}[language=typescript]
import type { HookHandler } from "../../src/hooks/hooks.js";

const handler: HookHandler = async (event) => {
  if (event.type !== "command" || event.action !== "new") {
    return;
  }

  console.log("[my-hook] Running!");
  // Your logic here
};

export default handler;
\end{lstlisting}


\subsubsection{5. 启用并测试}


\begin{lstlisting}[language=bash]
# Verify hook is discovered
openclaw hooks list

# Enable it
openclaw hooks enable my-hook

# Restart your gateway process (menu bar app restart on macOS, or restart your dev process)

# Trigger the event
# Send /new via your messaging channel
\end{lstlisting}


\subsection{配置}


\subsubsection{新配置格式（推荐）}


\begin{lstlisting}[language=json]
{
  "hooks": {
    "internal": {
      "enabled": true,
      "entries": {
        "session-memory": { "enabled": true },
        "command-logger": { "enabled": false }
      }
    }
  }
}
\end{lstlisting}


\subsubsection{每 Hook 配置}


Hooks 可以有自定义配置：

\begin{lstlisting}[language=json]
{
  "hooks": {
    "internal": {
      "enabled": true,
      "entries": {
        "my-hook": {
          "enabled": true,
          "env": {
            "MY_CUSTOM_VAR": "value"
          }
        }
      }
    }
  }
}
\end{lstlisting}


\subsubsection{额外目录}


从额外目录加载 hooks：

\begin{lstlisting}[language=json]
{
  "hooks": {
    "internal": {
      "enabled": true,
      "load": {
        "extraDirs": ["/path/to/more/hooks"]
      }
    }
  }
}
\end{lstlisting}


\subsubsection{遗留配置格式（仍然支持）}


旧配置格式仍然有效以保持向后兼容：

\begin{lstlisting}[language=json]
{
  "hooks": {
    "internal": {
      "enabled": true,
      "handlers": [
        {
          "event": "command:new",
          "module": "./hooks/handlers/my-handler.ts",
          "export": "default"
        }
      ]
    }
  }
}
\end{lstlisting}


\textbf{迁移}：对新 hooks 使用基于发现的新系统。遗留处理程序在基于目录的 hooks 之后加载。

\subsection{CLI 命令}


\subsubsection{列出 Hooks}


\begin{lstlisting}[language=bash]
# List all hooks
openclaw hooks list

# Show only eligible hooks
openclaw hooks list --eligible

# Verbose output (show missing requirements)
openclaw hooks list --verbose

# JSON output
openclaw hooks list --json
\end{lstlisting}


\subsubsection{Hook 信息}


\begin{lstlisting}[language=bash]
# Show detailed info about a hook
openclaw hooks info session-memory

# JSON output
openclaw hooks info session-memory --json
\end{lstlisting}


\subsubsection{检查资格}


\begin{lstlisting}[language=bash]
# Show eligibility summary
openclaw hooks check

# JSON output
openclaw hooks check --json
\end{lstlisting}


\subsubsection{启用/禁用}


\begin{lstlisting}[language=bash]
# Enable a hook
openclaw hooks enable session-memory

# Disable a hook
openclaw hooks disable command-logger
\end{lstlisting}


\subsection{捆绑的 Hooks}


\subsubsection{session-memory}


当你发出 \texttt{/new} 时将会话上下文保存到记忆。

\textbf{事件}：\texttt{command:new}

\textbf{要求}：必须配置 \texttt{workspace.dir}

\textbf{输出}：\texttt{/memory/YYYY-MM-DD-slug.md}（默认为 \texttt{\textasciitilde{}/.openclaw/workspace}）

\textbf{功能}：

\begin{enumerate}
  \item 使用预重置会话条目定位正确的记录
  \item 提取最后 15 行对话
  \item 使用 LLM 生成描述性文件名 slug
  \item 将会话元数据保存到带日期的记忆文件
\end{enumerate}


\textbf{示例输出}：

\begin{lstlisting}
# Session: 2026-01-16 14:30:00 UTC

- **Session Key**: agent:main:main
- **Session ID**: abc123def456
- **Source**: telegram
\end{lstlisting}


\textbf{文件名示例}：

\begin{itemize}
  \item \texttt{2026-01-16-vendor-pitch.md}
  \item \texttt{2026-01-16-api-design.md}
  \item \texttt{2026-01-16-1430.md}（如果 slug 生成失败则回退到时间戳）
\end{itemize}


\textbf{启用}：

\begin{lstlisting}[language=bash]
openclaw hooks enable session-memory
\end{lstlisting}


\subsubsection{command-logger}


将所有命令事件记录到集中审计文件。

\textbf{事件}：\texttt{command}

\textbf{要求}：无

\textbf{输出}：\texttt{\textasciitilde{}/.openclaw/logs/commands.log}

\textbf{功能}：

\begin{enumerate}
  \item 捕获事件详情（命令操作、时间戳、会话键、发送者 ID、来源）
  \item 以 JSONL 格式追加到日志文件
  \item 在后台静默运行
\end{enumerate}


\textbf{示例日志条目}：

\begin{lstlisting}
{"timestamp":"2026-01-16T14:30:00.000Z","action":"new","sessionKey":"agent:main:main","senderId":"+1234567890","source":"telegram"}
{"timestamp":"2026-01-16T15:45:22.000Z","action":"stop","sessionKey":"agent:main:main","senderId":"user@example.com","source":"whatsapp"}
\end{lstlisting}


\textbf{查看日志}：

\begin{lstlisting}[language=bash]
# View recent commands
tail -n 20 ~/.openclaw/logs/commands.log

# Pretty-print with jq
cat ~/.openclaw/logs/commands.log | jq .

# Filter by action
grep '"action":"new"' ~/.openclaw/logs/commands.log | jq .
\end{lstlisting}


\textbf{启用}：

\begin{lstlisting}[language=bash]
openclaw hooks enable command-logger
\end{lstlisting}


\subsubsection{boot-md}


当 Gateway 网关启动时运行 \texttt{BOOT.md}（在渠道启动之后）。
必须启用内部 hooks 才能运行。

\textbf{事件}：\texttt{gateway:startup}

\textbf{要求}：必须配置 \texttt{workspace.dir}

\textbf{功能}：

\begin{enumerate}
  \item 从你的工作区读取 \texttt{BOOT.md}
  \item 通过智能体运行器运行指令
  \item 通过 message 工具发送任何请求的出站消息
\end{enumerate}


\textbf{启用}：

\begin{lstlisting}[language=bash]
openclaw hooks enable boot-md
\end{lstlisting}


\subsection{最佳实践}


\subsubsection{保持处理程序快速}


Hooks 在命令处理期间运行。保持它们轻量：

\begin{lstlisting}[language=typescript]
// ✓ Good - async work, returns immediately
const handler: HookHandler = async (event) => {
  void processInBackground(event); // Fire and forget
};

// ✗ Bad - blocks command processing
const handler: HookHandler = async (event) => {
  await slowDatabaseQuery(event);
  await evenSlowerAPICall(event);
};
\end{lstlisting}


\subsubsection{优雅处理错误}


始终包装有风险的操作：

\begin{lstlisting}[language=typescript]
const handler: HookHandler = async (event) => {
  try {
    await riskyOperation(event);
  } catch (err) {
    console.error("[my-handler] Failed:", err instanceof Error ? err.message : String(err));
    // Don't throw - let other handlers run
  }
};
\end{lstlisting}


\subsubsection{尽早过滤事件}


如果事件不相关则尽早返回：

\begin{lstlisting}[language=typescript]
const handler: HookHandler = async (event) => {
  // Only handle 'new' commands
  if (event.type !== "command" || event.action !== "new") {
    return;
  }

  // Your logic here
};
\end{lstlisting}


\subsubsection{使用特定事件键}


尽可能在元数据中指定确切事件：

\begin{lstlisting}[language=yaml]
metadata: { "openclaw": { "events": ["command:new"] } } # Specific
\end{lstlisting}


而不是：

\begin{lstlisting}[language=yaml]
metadata: { "openclaw": { "events": ["command"] } } # General - more overhead
\end{lstlisting}


\subsection{调试}


\subsubsection{启用 Hook 日志}


Gateway 网关在启动时记录 hook 加载：

\begin{lstlisting}
Registered hook: session-memory -> command:new
Registered hook: command-logger -> command
Registered hook: boot-md -> gateway:startup
\end{lstlisting}


\subsubsection{检查发现}


列出所有发现的 hooks：

\begin{lstlisting}[language=bash]
openclaw hooks list --verbose
\end{lstlisting}


\subsubsection{检查注册}


在你的处理程序中，记录它被调用的时间：

\begin{lstlisting}[language=typescript]
const handler: HookHandler = async (event) => {
  console.log("[my-handler] Triggered:", event.type, event.action);
  // Your logic
};
\end{lstlisting}


\subsubsection{验证资格}


检查为什么 hook 不符合条件：

\begin{lstlisting}[language=bash]
openclaw hooks info my-hook
\end{lstlisting}


在输出中查找缺失的要求。

\subsection{测试}


\subsubsection{Gateway 网关日志}


监控 Gateway 网关日志以查看 hook 执行：

\begin{lstlisting}[language=bash]
# macOS
./scripts/clawlog.sh -f

# Other platforms
tail -f ~/.openclaw/gateway.log
\end{lstlisting}


\subsubsection{直接测试 Hooks}


隔离测试你的处理程序：

\begin{lstlisting}[language=typescript]
import { test } from "vitest";
import { createHookEvent } from "./src/hooks/hooks.js";
import myHandler from "./hooks/my-hook/handler.js";

test("my handler works", async () => {
  const event = createHookEvent("command", "new", "test-session", {
    foo: "bar",
  });

  await myHandler(event);

  // Assert side effects
});
\end{lstlisting}


\subsection{架构}


\subsubsection{核心组件}


\begin{itemize}
  \item \textbf{\texttt{src/hooks/types.ts}}：类型定义
  \item \textbf{\texttt{src/hooks/workspace.ts}}：目录扫描和加载
  \item \textbf{\texttt{src/hooks/frontmatter.ts}}：HOOK.md 元数据解析
  \item \textbf{\texttt{src/hooks/config.ts}}：资格检查
  \item \textbf{\texttt{src/hooks/hooks-status.ts}}：状态报告
  \item \textbf{\texttt{src/hooks/loader.ts}}：动态模块加载器
  \item \textbf{\texttt{src/cli/hooks-cli.ts}}：CLI 命令
  \item \textbf{\texttt{src/gateway/server-startup.ts}}：在 Gateway 网关启动时加载 hooks
  \item \textbf{\texttt{src/auto-reply/reply/commands-core.ts}}：触发命令事件
\end{itemize}


\subsubsection{发现流程}


\begin{lstlisting}
Gateway 网关启动
    ↓
扫描目录（工作区 → 托管 → 捆绑）
    ↓
解析 HOOK.md 文件
    ↓
检查资格（bins、env、config、os）
    ↓
从符合条件的 hooks 加载处理程序
    ↓
为事件注册处理程序
\end{lstlisting}


\subsubsection{事件流程}


\begin{lstlisting}
用户发送 /new
    ↓
命令验证
    ↓
创建 hook 事件
    ↓
触发 hook（所有注册的处理程序）
    ↓
命令处理继续
    ↓
会话重置
\end{lstlisting}


\subsection{故障排除}


\subsubsection{Hook 未被发现}


\begin{enumerate}
  \item 检查目录结构：
\end{enumerate}


\begin{lstlisting}[language=bash]
   ls -la ~/.openclaw/hooks/my-hook/
   # Should show: HOOK.md, handler.ts
\end{lstlisting}


\begin{enumerate}
  \item 验证 HOOK.md 格式：
\end{enumerate}


\begin{lstlisting}[language=bash]
   cat ~/.openclaw/hooks/my-hook/HOOK.md
   # Should have YAML frontmatter with name and metadata
\end{lstlisting}


\begin{enumerate}
  \item 列出所有发现的 hooks：
\begin{lstlisting}[language=bash]
   openclaw hooks list
\end{lstlisting}

\end{enumerate}


\subsubsection{Hook 不符合条件}


检查要求：

\begin{lstlisting}[language=bash]
openclaw hooks info my-hook
\end{lstlisting}


查找缺失的：

\begin{itemize}
  \item 二进制文件（检查 PATH）
  \item 环境变量
  \item 配置值
  \item 操作系统兼容性
\end{itemize}


\subsubsection{Hook 未执行}


\begin{enumerate}
  \item 验证 hook 已启用：
\end{enumerate}


\begin{lstlisting}[language=bash]
   openclaw hooks list
   # Should show ✓ next to enabled hooks
\end{lstlisting}


\begin{enumerate}
  \item 重启你的 Gateway 网关进程以重新加载 hooks。
\end{enumerate}


\begin{enumerate}
  \item 检查 Gateway 网关日志中的错误：
\begin{lstlisting}[language=bash]
   ./scripts/clawlog.sh | grep hook
\end{lstlisting}

\end{enumerate}


\subsubsection{处理程序错误}


检查 TypeScript/import 错误：

\begin{lstlisting}[language=bash]
# Test import directly
node -e "import('./path/to/handler.ts').then(console.log)"
\end{lstlisting}


\subsection{迁移指南}


\subsubsection{从遗留配置到发现}


\textbf{之前}：

\begin{lstlisting}[language=json]
{
  "hooks": {
    "internal": {
      "enabled": true,
      "handlers": [
        {
          "event": "command:new",
          "module": "./hooks/handlers/my-handler.ts"
        }
      ]
    }
  }
}
\end{lstlisting}


\textbf{之后}：

\begin{enumerate}
  \item 创建 hook 目录：
\end{enumerate}


\begin{lstlisting}[language=bash]
   mkdir -p ~/.openclaw/hooks/my-hook
   mv ./hooks/handlers/my-handler.ts ~/.openclaw/hooks/my-hook/handler.ts
\end{lstlisting}


\begin{enumerate}
  \item 创建 HOOK.md：
\end{enumerate}


\begin{lstlisting}
   ---
   name: my-hook
   description: "My custom hook"
   metadata: { "openclaw": { "emoji": "🎯", "events": ["command:new"] } }
   ---

   # My Hook

   Does something useful.
\end{lstlisting}


\begin{enumerate}
  \item 更新配置：
\end{enumerate}


\begin{lstlisting}[language=json]
   {
     "hooks": {
       "internal": {
         "enabled": true,
         "entries": {
           "my-hook": { "enabled": true }
         }
       }
     }
   }
\end{lstlisting}


\begin{enumerate}
  \item 验证并重启你的 Gateway 网关进程：
\begin{lstlisting}[language=bash]
   openclaw hooks list
   # Should show: 🎯 my-hook ✓
\end{lstlisting}

\end{enumerate}


\textbf{迁移的好处}：

\begin{itemize}
  \item 自动发现
  \item CLI 管理
  \item 资格检查
  \item 更好的文档
  \item 一致的结构
\end{itemize}


\subsection{另请参阅}


\begin{itemize}
  \item \href{/cli/hooks}{CLI 参考：hooks}
  \item \href{https://github.com/openclaw/openclaw/tree/main/src/hooks/bundled}{捆绑 Hooks README}
  \item \href{/automation/webhook}{Webhook Hooks}
  \item \href{/gateway/configuration\#hooks}{配置}
\end{itemize}



\clearpage

% === 源文件: automation/poll.md ===
\section{投票}


\subsection{支持的渠道}


\begin{itemize}
  \item WhatsApp（Web 渠道）
  \item Discord
  \item MS Teams（Adaptive Cards）
\end{itemize}


\subsection{CLI}


\begin{lstlisting}[language=bash]
# WhatsApp
openclaw message poll --target +15555550123 \
  --poll-question "Lunch today?" --poll-option "Yes" --poll-option "No" --poll-option "Maybe"
openclaw message poll --target 123456789@g.us \
  --poll-question "Meeting time?" --poll-option "10am" --poll-option "2pm" --poll-option "4pm" --poll-multi

# Discord
openclaw message poll --channel discord --target channel:123456789 \
  --poll-question "Snack?" --poll-option "Pizza" --poll-option "Sushi"
openclaw message poll --channel discord --target channel:123456789 \
  --poll-question "Plan?" --poll-option "A" --poll-option "B" --poll-duration-hours 48

# MS Teams
openclaw message poll --channel msteams --target conversation:19:abc@thread.tacv2 \
  --poll-question "Lunch?" --poll-option "Pizza" --poll-option "Sushi"
\end{lstlisting}


选项：

\begin{itemize}
  \item \texttt{--channel}：\texttt{whatsapp}（默认）、\texttt{discord} 或 \texttt{msteams}
  \item \texttt{--poll-multi}：允许选择多个选项
  \item \texttt{--poll-duration-hours}：仅限 Discord（省略时默认为 24）
\end{itemize}


\subsection{Gateway 网关 RPC}


方法：\texttt{poll}

参数：

\begin{itemize}
  \item \texttt{to}（字符串，必需）
  \item \texttt{question}（字符串，必需）
  \item \texttt{options}（字符串数组，必需）
  \item \texttt{maxSelections}（数字，可选）
  \item \texttt{durationHours}（数字，可选）
  \item \texttt{channel}（字符串，可选，默认：\texttt{whatsapp}）
  \item \texttt{idempotencyKey}（字符串，必需）
\end{itemize}


\subsection{渠道差异}


\begin{itemize}
  \item WhatsApp：2-12 个选项，\texttt{maxSelections} 必须在选项数量范围内，忽略 \texttt{durationHours}。
  \item Discord：2-10 个选项，\texttt{durationHours} 限制在 1-768 小时之间（默认 24）。\texttt{maxSelections > 1} 启用多选；Discord 不支持严格的选择数量限制。
  \item MS Teams：Adaptive Card 投票（由 OpenClaw 管理）。无原生投票 API；\texttt{durationHours} 被忽略。
\end{itemize}


\subsection{智能体工具（Message）}


使用 \texttt{message} 工具的 \texttt{poll} 操作（\texttt{to}、\texttt{pollQuestion}、\texttt{pollOption}，可选 \texttt{pollMulti}、\texttt{pollDurationHours}、\texttt{channel}）。

注意：Discord 没有"恰好选择 N 个"模式；\texttt{pollMulti} 映射为多选。
Teams 投票以 Adaptive Cards 形式渲染，需要 Gateway 网关保持在线
以将投票记录到 \texttt{\textasciitilde{}/.openclaw/msteams-polls.json}。


\clearpage

% === 源文件: automation/troubleshooting.md ===
\section{自动化故障排查}


该页面是英文文档的中文占位版本，完整内容请先参考英文版：\href{/automation/troubleshooting}{Automation Troubleshooting}。


\clearpage

% === 源文件: automation/webhook.md ===
\section{Webhooks}


Gateway 网关可以暴露一个小型 HTTP webhook 端点用于外部触发。

\subsection{启用}


\begin{lstlisting}[language=json]
{
  hooks: {
    enabled: true,
    token: "shared-secret",
    path: "/hooks",
  },
}
\end{lstlisting}


注意事项：

\begin{itemize}
  \item 当 \texttt{hooks.enabled=true} 时，\texttt{hooks.token} 为必填项。
  \item \texttt{hooks.path} 默认为 \texttt{/hooks}。
\end{itemize}


\subsection{认证}


每个请求必须包含 hook 令牌。推荐使用请求头：

\begin{itemize}
  \item \texttt{Authorization: Bearer }（推荐）
  \item \texttt{x-openclaw-token: }
  \item \texttt{?token=}（已弃用；会记录警告日志，将在未来的主要版本中移除）
\end{itemize}


\subsection{端点}


\subsubsection{POST /hooks/wake}


请求体：

\begin{lstlisting}[language=json]
{ "text": "System line", "mode": "now" }
\end{lstlisting}


\begin{itemize}
  \item \texttt{text} \textbf{必填}（字符串）：事件描述（例如"收到新邮件"）。
  \item \texttt{mode} 可选（\texttt{now} | \texttt{next-heartbeat}）：是否立即触发心跳（默认 \texttt{now}）或等待下一次定期检查。
\end{itemize}


效果：

\begin{itemize}
  \item 为\textbf{主}会话加入一个系统事件队列
  \item 如果 \texttt{mode=now}，则立即触发心跳
\end{itemize}


\subsubsection{POST /hooks/agent}


请求体：

\begin{lstlisting}[language=json]
{
  "message": "Run this",
  "name": "Email",
  "sessionKey": "hook:email:msg-123",
  "wakeMode": "now",
  "deliver": true,
  "channel": "last",
  "to": "+15551234567",
  "model": "openai/gpt-5.2-mini",
  "thinking": "low",
  "timeoutSeconds": 120
}
\end{lstlisting}


\begin{itemize}
  \item \texttt{message} \textbf{必填}（字符串）：智能体要处理的提示或消息。
  \item \texttt{name} 可选（字符串）：hook 的可读名称（例如"GitHub"），用作会话摘要的前缀。
  \item \texttt{sessionKey} 可选（字符串）：用于标识智能体会话的键。默认为随机的 \texttt{hook:}。使用一致的键可以在 hook 上下文中进行多轮对话。
  \item \texttt{wakeMode} 可选（\texttt{now} | \texttt{next-heartbeat}）：是否立即触发心跳（默认 \texttt{now}）或等待下一次定期检查。
  \item \texttt{deliver} 可选（布尔值）：如果为 \texttt{true}，智能体的响应将发送到消息渠道。默认为 \texttt{true}。仅为心跳确认的响应会自动跳过。
  \item \texttt{channel} 可选（字符串）：用于投递的消息渠道。可选值：\texttt{last}、\texttt{whatsapp}、\texttt{telegram}、\texttt{discord}、\texttt{slack}、\texttt{mattermost}（插件）、\texttt{signal}、\texttt{imessage}、\texttt{msteams}。默认为 \texttt{last}。
  \item \texttt{to} 可选（字符串）：渠道的接收者标识符（例如 WhatsApp/Signal 的电话号码、Telegram 的聊天 ID、Discord/Slack/Mattermost（插件）的频道 ID、MS Teams 的会话 ID）。默认为主会话中的最后一个接收者。
  \item \texttt{model} 可选（字符串）：模型覆盖（例如 \texttt{anthropic/claude-3-5-sonnet} 或别名）。如果有限制，必须在允许的模型列表中。
  \item \texttt{thinking} 可选（字符串）：思考级别覆盖（例如 \texttt{low}、\texttt{medium}、\texttt{high}）。
  \item \texttt{timeoutSeconds} 可选（数字）：智能体运行的最大持续时间（秒）。
\end{itemize}


效果：

\begin{itemize}
  \item 运行一个\textbf{隔离的}智能体回合（独立的会话键）
  \item 始终在\textbf{主}会话中发布摘要
  \item 如果 \texttt{wakeMode=now}，则立即触发心跳
\end{itemize}


\subsubsection{POST /hooks/（映射）}


自定义 hook 名称通过 \texttt{hooks.mappings} 解析（见配置）。映射可以将任意请求体转换为 \texttt{wake} 或 \texttt{agent} 操作，支持可选的模板或代码转换。

映射选项（摘要）：

\begin{itemize}
  \item \texttt{hooks.presets: ["gmail"]} 启用内置的 Gmail 映射。
  \item \texttt{hooks.mappings} 允许你在配置中定义 \texttt{match}、\texttt{action} 和模板。
  \item \texttt{hooks.transformsDir} + \texttt{transform.module} 加载 JS/TS 模块用于自定义逻辑。
  \item 使用 \texttt{match.source} 保持通用的接收端点（基于请求体的路由）。
  \item TS 转换需要 TS 加载器（例如 \texttt{bun} 或 \texttt{tsx}）或运行时预编译的 \texttt{.js}。
  \item 在映射上设置 \texttt{deliver: true} + \texttt{channel}/\texttt{to} 可将回复路由到聊天界面（\texttt{channel} 默认为 \texttt{last}，回退到 WhatsApp）。
  \item \texttt{allowUnsafeExternalContent: true} 禁用该 hook 的外部内容安全包装（危险；仅用于受信任的内部来源）。
  \item \texttt{openclaw webhooks gmail setup} 为 \texttt{openclaw webhooks gmail run} 写入 \texttt{hooks.gmail} 配置。完整的 Gmail 监听流程请参阅 \href{/automation/gmail-pubsub}{Gmail Pub/Sub}。
\end{itemize}


\subsection{响应}


\begin{itemize}
  \item \texttt{200} 用于 \texttt{/hooks/wake}
  \item \texttt{202} 用于 \texttt{/hooks/agent}（异步运行已启动）
  \item \texttt{401} 认证失败
  \item \texttt{400} 请求体无效
  \item \texttt{413} 请求体过大
\end{itemize}


\subsection{示例}


\begin{lstlisting}[language=bash]
curl -X POST http://127.0.0.1:18789/hooks/wake \
  -H 'Authorization: Bearer SECRET' \
  -H 'Content-Type: application/json' \
  -d '{"text":"New email received","mode":"now"}'
\end{lstlisting}


\begin{lstlisting}[language=bash]
curl -X POST http://127.0.0.1:18789/hooks/agent \
  -H 'x-openclaw-token: SECRET' \
  -H 'Content-Type: application/json' \
  -d '{"message":"Summarize inbox","name":"Email","wakeMode":"next-heartbeat"}'
\end{lstlisting}


\subsubsection{使用不同的模型}


在智能体请求体（或映射）中添加 \texttt{model} 以覆盖该次运行的模型：

\begin{lstlisting}[language=bash]
curl -X POST http://127.0.0.1:18789/hooks/agent \
  -H 'x-openclaw-token: SECRET' \
  -H 'Content-Type: application/json' \
  -d '{"message":"Summarize inbox","name":"Email","model":"openai/gpt-5.2-mini"}'
\end{lstlisting}


如果你启用了 \texttt{agents.defaults.models} 限制，请确保覆盖的模型包含在其中。

\begin{lstlisting}[language=bash]
curl -X POST http://127.0.0.1:18789/hooks/gmail \
  -H 'Authorization: Bearer SECRET' \
  -H 'Content-Type: application/json' \
  -d '{"source":"gmail","messages":[{"from":"Ada","subject":"Hello","snippet":"Hi"}]}'
\end{lstlisting}


\subsection{安全}


\begin{itemize}
  \item 将 hook 端点保持在 loopback、tailnet 或受信任的反向代理之后。
  \item 使用专用的 hook 令牌；不要复用 Gateway 网关认证令牌。
  \item 避免在 webhook 日志中包含敏感的原始请求体。
  \item Hook 请求体默认被视为不受信任并使用安全边界包装。如果你必须为特定 hook 禁用此功能，请在该 hook 的映射中设置 \texttt{allowUnsafeExternalContent: true}（危险）。
\end{itemize}



\clearpage
